\section{总结与延伸}

课程最后从四个案例上升到 foundation biomedical AI。讲者的核心判断不是“是否”会发生变化，而是“何时、以何种方式”发生。要理解这句话，必须把前文重新压缩成一条完整链：临床模型处理人类知识与沟通，蛋白质模型连接 sequence 与 function language，DeepConsensus 提高观测质量，Enformer解释调控上下文；未来系统的机会在于让这些层互相提供证据，而不是把所有数据塞进一个聊天界面。

\subsection{从四个模型提炼出的共同框架}

未来章节页先标记讨论层次的变化，随后提出“when and how”。这不是时间预测，而是工程与科学问题：哪些数据能够合法、安全地联合？哪些输出可以由自动 metric 验证，哪些必须交给 clinician 或 wet lab？多模态 foundation model 的规模收益，如何与 provenance、calibration 和 selective prediction 一起设计？

\lecturefigure{slide-79-section-future.jpg}{章节切换：Biomedical AI 的未来}{官方录像恢复幻灯片 V079；Stanford CS25 Lecture 17}

\lecturefigure{slide-80-future-how-not-if.jpg}{未来问题：不是会不会，而是何时、如何发生}{官方录像恢复幻灯片 V080；Stanford CS25 Lecture 17}

\begin{longtable}{@{}L{0.19\textwidth}L{0.25\textwidth}L{0.25\textwidth}L{0.19\textwidth}@{}}
\toprule
案例 & 输入与挑战 & 模型机制 & 外部验证 \\
\midrule
Med-PaLM & 医学问题；知识、沟通与风险 & Scaling、instruction tuning、prompt tuning、abstention & 医生与普通用户多维评价 \\
Performer / Protein LM & 超长氨基酸序列 & kernel attention 与表示学习 & 结构、数据库与实验 \\
ProtNLM & 蛋白质序列到自由文本 & encoder--decoder、多任务 T5 & 注释 provenance、专家与实验 \\
DeepConsensus & 对齐的 noisy subreads & gap-aware attention、alignment loss、quality head & base/read accuracy 与系统吞吐 \\
Enformer & 长 DNA 到 genomic tracks & convolution + long-range attention & 实验轨迹、表达与 variant assay \\
\bottomrule
\end{longtable}

\begin{importantbox}{统一原则：模型输出必须接回对应世界}
文本答案接回临床共识与使用者，蛋白质描述接回数据库和实验，测序纠错接回 read 与 pipeline，调控预测接回 assay 与表型。Foundation model 可以共享表示和数据，但不能共享一个模糊的“看起来合理”评价标准。
\end{importantbox}

这张总表还说明，foundation model 的共享部分与专用部分应被分别设计。共享 encoder 可以学习跨模态表示、统一检索和条件生成；专用 output head、loss、calibration 与 reviewer 则负责把输出约束到临床、蛋白质或基因组语义。若为了“统一”删除专用验证，系统会在界面上更简洁，却在证据上更脆弱。相反，模块化并不妨碍端到端训练，只要求每个关键接口都能单独审计：输入来自哪里、表示怎样传递、输出由谁确认、失败如何升级。

\subsection{Foundation biomedical AI：整合而不是抹平差异}

Foundation biomedical AI 页用多模态圆盘表达整合愿景。下一页并列著名 physician 与 scientist，讲者借此说明未来 AI 不必把 clinical application 与 biological application 分开：临床观察可以提出机制问题，分子模型可以生成候选，实验结果再更新临床知识。这里的“foundation”更像跨层接口，而不是单纯更大参数量。

\lecturefigure{slide-81-foundation-biomedical-ai.jpg}{Foundation biomedical AI 的多模态整合愿景}{官方录像恢复幻灯片 V081；Stanford CS25 Lecture 17}

\lecturefigure{slide-82-biomedical-ai-scientists.jpg}{医生与科学家的角色可以在 AI 系统中被连接}{官方录像恢复幻灯片 V082；Stanford CS25 Lecture 17}

\teachervoice{讲者说 clinical applications 与 biological applications 不必永远分离；把它们连接起来，才可能产生新 insight、加速 biomedical research，并最终回到疾病与健康问题。这个判断应被写成研究方向，而不是声称 2023 年已经有统一闭环。}

\begin{warningbox}{跨层整合最容易发生的三种混淆}
第一，把临床相关性当成分子因果；第二，把模型 confidence 当成实验重复性；第三，把同一患者或数据库泄漏到训练和测试当成泛化。越是 foundation 化，数据 provenance、时间切分、隐私与独立验证越重要。
\end{warningbox}

\subsection{AI 会在哪个领域获得第一枚 Nobel Prize？}

倒数第二页把问题留给学生：AI 最先在哪个领域促成 Nobel-level discovery？这不是要求猜奖项，而是在追问“什么证据才算科学突破”。如果 AI 只是压缩已有文献，贡献属于工具；如果它提出新机制、指导实验并经独立复现，才接近发现。生物医学的特殊之处在于它同时拥有巨大未解空间、可测量实验和直接人类价值，也拥有最高的伦理与验证门槛。

\lecturefigure{slide-83-ai-nobel-question.jpg}{开放问题：AI 会在哪个领域获得第一枚 Nobel Prize？}{官方录像恢复幻灯片 V083；Stanford CS25 Lecture 17}

\lecturefigure{slide-84-thank-you.jpg}{课程结束与协作者致谢}{官方录像恢复幻灯片 V084；Stanford CS25 Lecture 17}

\teachervoice{讲者把问题留在课堂，没有给出年份和确定答案；随后表示自己希望多数人会想到 medicine。真正值得保留的是评价尺度：AI 若要在医学产生 Nobel-level 影响，必须从预测走到可复现发现，再从发现走到健康收益。}

\subsection{给工程实践的七条检查清单}

\begin{enumerate}
  \item \textbf{先定义序列单位。} token、event、amino acid 与 nucleotide 的位置语义不能混用。
  \item \textbf{先写外部验证，再选模型。} 如果输出无法接回 clinician、database、assay 或 pipeline，benchmark 很难代表真实价值。
  \item \textbf{把长上下文成本显式化。} 记录 sequence length、attention 复杂度、截断策略与可能丢失的远程关系。
  \item \textbf{把不确定性做成接口。} selective prediction、quality score 与 calibration 必须服务后续升级路径，而不是只画一条曲线。
  \item \textbf{把多维风险保留下来。} accuracy、missing content、harm、bias 与 helpfulness 不应过早压成一个分数。
  \item \textbf{区分预测、关联和因果。} attention map、variant score 与生成注释都可以提出假设，但不能替代实验。
  \item \textbf{保留历史边界。} 用课堂当时的模型和证据回答课堂问题，后来的进展应放入独立更新，而不是改写原讲义的因果链。
\end{enumerate}

\begin{importantbox}{最终压缩}
Biomedical Transformer 的核心不是把医学、生物学都“语言化”，而是用统一的上下文建模工具连接不同序列，同时为每个领域保留专属表示、损失、校准和验证。真正的 foundation system 不是一个无所不答的模型，而是一组能在多层证据之间传递信息、知道何时拒答、并把关键判断交还给医生和实验的可靠接口。
\end{importantbox}

\subsection{拓展阅读}

\begin{itemize}
  \item Singhal et al., \href{https://www.nature.com/articles/s41586-023-06291-2}{\emph{Large Language Models Encode Clinical Knowledge}（Nature）}
  \item Choromanski et al., \href{https://arxiv.org/abs/2009.14794}{\emph{Rethinking Attention with Performers}（arXiv）}
  \item Google Research, \href{https://github.com/google-research/google-research/tree/master/protnlm}{\emph{ProtNLM} 官方代码与说明}
  \item Baid et al., \href{https://research.google/pubs/deepconsensus-improves-the-accuracy-of-sequences-with-a-gap-aware-sequence-transformer/}{\emph{DeepConsensus}（Google Research）}
  \item Avsec et al., \href{https://www.nature.com/articles/s41592-021-01252-x}{\emph{Effective Gene Expression Prediction from Sequence by Integrating Long-range Interactions}（Nature Methods）}
\end{itemize}

这五项工作分别代表临床评价、长序列 attention、科学语言接口、测序基础设施和调控基因组学。建议阅读时统一使用本讲的五个问题：输入怎样表示、上下文为何困难、监督从哪里来、输出如何验证、结论不能推出什么。这样才能把论文结果转化为可迁移的研究方法，而不是模型名词清单。

\end{document}
